% Sampul depan, versi formal. Dipanggil dari main.tex lewat \versiSampul.
\begin{tikzpicture}[remember picture,overlay]
  \fill[merah] (current page.south west) rectangle (current page.north east);
  % Pasangannya "{" ada di sampul belakang
  \node[text=marun!55!merah,font=\judul\fontsize{470}{470}\selectfont,inner sep=0pt]
    at ([xshift=-14mm,yshift=-100mm]current page.north east) {\}};
  % Blok bawah
  \fill[marun] (current page.south west) rectangle ([yshift=104mm]current page.south east);
  \fill[emas] ([yshift=104mm]current page.south west) rectangle ([yshift=105.2mm]current page.south east);
  % Logo dan institusi
  \node[anchor=north west,inner sep=0pt] (logo) at ([xshift=22mm,yshift=-20mm]current page.north west) {\includegraphics[height=20mm]{\logoFile}};
  \node[anchor=south west,text=white,font=\capsfont\small,inner sep=0pt] at ([xshift=6mm,yshift=1mm]logo.east) {\universitas};
  \node[anchor=north west,text=emasmuda,font=\capsfont\scriptsize,inner sep=0pt] at ([xshift=6mm,yshift=-2mm]logo.east) {\fakultas};
  % Judul
  \fill[emas] ([xshift=22mm,yshift=-70mm]current page.north west) rectangle ++(14mm,-1.2pt);
  \node[anchor=west,text=emasmuda,font=\capsfont\small,inner sep=0pt] at ([xshift=40mm,yshift=-70.2mm]current page.north west) {\jenisBuku};
  \node[anchor=north west,text=white,font=\judul\fontsize{60}{58}\selectfont,text width=150mm,align=left,inner sep=0pt]
    (judul) at ([xshift=22mm,yshift=-80mm]current page.north west) {\judulBuku};
  \node[anchor=north west,text=emasmuda,font=\itshape\fontsize{20}{25}\selectfont,text width=150mm,align=left,inner sep=0pt]
    (sub) at ([yshift=-8mm]judul.south west) {\subJudul};
  % Diagram kelas
  \node[kelas,anchor=north] (induk) at ([yshift=94mm]current page.south) {\kelas{Civitas}{- nama : str\\- email : str}{+ perkenalan() : str\\+ kontak() : str}};
  \node[kelas,anchor=north] (anak1) at ([xshift=-32mm,yshift=-18mm]induk.south) {\kelas{Mahasiswa}{- nim : str}{+ isi\_krs(matkul) : str}};
  \node[kelas,anchor=north] (anak2) at ([xshift=32mm,yshift=-18mm]induk.south) {\kelas{Dosen}{- nip : str}{+ beri\_nilai(mhs) : str}};
  \coordinate (cabang) at ([yshift=-9mm]induk.south);
  \draw[emas,line width=0.9pt] (anak1.north) |- (cabang) -| (anak2.north);
  \draw[turunan,emas] (cabang) -- (induk.south);
  % Kaki sampul
  \node[anchor=south west,text=emasmuda,font=\capsfont\scriptsize] at ([xshift=22mm,yshift=10mm]current page.south west)
    {\prodi\quad\textbullet\quad \tahun};
\end{tikzpicture}
\clearpage

